\section*{Розділ V. Спеціальні процедури}
\addcontentsline{toc}{section}{Розділ V. Спеціальні процедури}

\subsection*{5.1. Процедура затвердження положень та регламентів}
\addcontentsline{toc}{subsection}{5.1. Процедура затвердження положень та регламентів}
    5.1.1. Проєкти положень про органи студентського самоврядування та регламентів їх роботи розглядаються КСУ у спеціальному порядку з урахуванням їх системного значення для студентського самоврядування.

    5.1.2. Процедура затвердження включає такі етапи:

        \begin{enumerate}[label=\alph*)]
            \item Внесення проєкту документа на розгляд КСУ суб'єктом, що має право законодавчої ініціативи;
            \item Направлення проєкту до Студентської уповноваженої делегації (СУД) для отримання висновку щодо відповідності Положенню про ОСС та чинним нормативним актам студентського самоврядування;
            \item Розгляд висновку СУД та доопрацювання проєкту (за потреби);
            \item Обговорення проєкту на засіданні КСУ;
            \item Голосування та прийняття рішення.
        \end{enumerate}

    5.1.3. СУД надає висновок щодо проєкту протягом 10 робочих днів з дня його отримання. Висновок може містити зауваження щодо відповідності проєкту чинним нормам та пропозиції щодо його вдосконалення.

    5.1.4. У разі надання СУД негативного висновку, проєкт може бути винесений на голосування КСУ лише після його доопрацювання відповідно до зауважень СУД або за рішенням КСУ про розгляд без урахування зауважень СУД.

\subsection*{5.2. Процедура заслуховування звітів}
\addcontentsline{toc}{subsection}{5.2. Процедура заслуховування звітів}
    5.2.1. Звіти органів студентського самоврядування заслуховуються КСУ відповідно до пункту 3.2 Положення про ОСС за стандартною процедурою.

    5.2.2. Процедура заслуховування звіту включає:

        \begin{enumerate}[label=\alph*)]
            \item Доповідь керівника органу про діяльність за звітний період;
            \item Відповіді на запитання делегатів КСУ;
            \item Обговорення результатів діяльності;
            \item Голосування за затвердження звіту.
        \end{enumerate}

    5.2.3. Звіт вважається затвердженим, якщо за його затвердження проголосувала проста більшість від загального складу делегатів КСУ.

    5.2.4. У разі незатвердження звіту, КСУ може прийняти рішення про надання органу додаткового часу для усунення недоліків та повторного подання звіту.

\subsection*{5.3. Процедура висловлення недовіри}
\addcontentsline{toc}{subsection}{5.3. Процедура висловлення недовіри}
    5.3.1. Питання про висловлення недовіри керівникам та членам органів студентського самоврядування розглядається КСУ відповідно до пункту 3.2 Положення про ОСС з дотриманням спеціальної процедури.

    5.3.2. Ініціатива про висловлення недовіри може бути внесена:

        \begin{enumerate}[label=\alph*)]
            \item Студентською уповноваженою делегацією (СУД) за результатами розслідування порушень;
            \item Групою делегатів КСУ кількістю не менше п'яти осіб;
            \item Органом студентського самоврядування щодо свого керівника або члена.
        \end{enumerate}

    5.3.3. До розгляду питання про недовіру може бути залучена СУД для проведення попереднього розслідування обставин, що стали підставою для ініціативи про недовіру.

    5.3.4. Процедура розгляду включає:

        \begin{enumerate}[label=\alph*)]
            \item Заслуховування доповіді ініціатора про підстави для висловлення недовіри;
            \item Заслуховування висновку СУД (у разі проведення розслідування);
            \item Надання особі, щодо якої розглядається питання, права на захист та пояснення;
            \item Обговорення питання делегатами КСУ;
            \item Голосування (поіменне або таємне за вимогою особи).
        \end{enumerate}

    5.3.5. Рішення про висловлення недовіри приймається не менше ніж двома третинами голосів від загального складу делегатів КСУ та має наслідком автоматичне припинення повноважень відповідної особи.

\subsection*{5.4. Процедура затвердження кошторису ОСС}
\addcontentsline{toc}{subsection}{5.4. Процедура затвердження кошторису ОСС}
    5.4.1. Єдиний кошторис (бюджет) ОСС Університету затверджується КСУ відповідно до пункту 3.2 Положення про ОСС з урахуванням особливостей погодження Студентською уповноваженою делегацією.

    5.4.2. Процедура розгляду та затвердження кошторису з медіацією між КСУ та СУД включає:

        \begin{enumerate}[label=\alph*)]
            \item Подання проєкту кошторису Фінансовим комітетом КСУ до СУД для попереднього погодження не пізніше ніж за 15 робочих днів до засідання КСУ;
            \item У разі незгоди СУД з окремими положеннями проєкту -- проведення консультацій між Фінансовим комітетом КСУ, СУД та Спікером КСУ для узгодження позицій;
            \item Доопрацювання проєкту кошторису за результатами консультацій (за можливості);
            \item Внесення проєкту кошторису на розгляд КСУ з висновком СУД щодо погодження;
            \item Обговорення проєкту на засіданні КСУ;
            \item Голосування за затвердження.
        \end{enumerate}

    5.4.3. У разі погодження СУД проєкт кошторису затверджується простою більшістю від загального складу делегатів КСУ.

    5.4.4. У разі відсутності погодження СУД кошторис може бути затверджений лише за умови, якщо за нього проголосувало не менше двох третин (2/3) від загального складу делегатів КСУ.

    5.4.5. Консультації між органами ОСС з питань кошторису проводяться з метою досягнення компромісу та можуть включати:

        \begin{enumerate}[label=\alph*)]
            \item Зустрічі представників Фінансового комітету КСУ та СУД;
            \item Участь Спікера КСУ як медіатора;
            \item Залучення експертів з фінансових питань;
            \item Розробка альтернативних варіантів розподілу коштів.
        \end{enumerate}

\subsection*{5.5. Процедура розгляду апеляцій на рішення СУД}
\addcontentsline{toc}{subsection}{5.5. Процедура розгляду апеляцій на рішення СУД}
    5.5.1. Апеляції на рішення СУД розглядаються КСУ відповідно до Розділу VI.4 Положення про СУД з дотриманням повної процедури, передбаченої цим Положенням.

    5.5.2. Розгляд апеляції здійснюється у такому порядку:

        \begin{enumerate}[label=\alph*)]
            \item Створення тимчасової апеляційної комісії КСУ у складі 3-5 осіб з числа делегатів, що не мають конфлікту інтересів у справі;
            \item Вивчення комісією матеріалів справи та рішення СУД;
            \item Заслуховування комісією пояснень апелянта та представника СУД;
            \item Підготовка комісією висновку та проєкту рішення КСУ;
            \item Розгляд висновку комісії на пленарному засіданні КСУ;
            \item Прийняття КСУ остаточного рішення за апеляцією.
        \end{enumerate}

    5.5.3. За результатами розгляду апеляції КСУ більшістю голосів від загального складу делегатів може прийняти одне з таких рішень:

        \begin{enumerate}[label=\alph*)]
            \item Залишити рішення СУД без змін, а апеляцію без задоволення;
            \item Скасувати рішення СУД повністю або частково;
            \item Змінити рішення СУД;
            \item Направити справу на новий розгляд до СУД з наданням відповідних вказівок.
        \end{enumerate}

    5.5.4. Рішення КСУ за результатами розгляду апеляції є остаточним в системі студентського самоврядування та не підлягає оскарженню.

\subsection*{5.6. Процедура призначення виконувачів обов'язків}
\addcontentsline{toc}{subsection}{5.6. Процедура призначення виконувачів обов'язків}
    5.6.1. КСУ розглядає питання про затвердження або відхилення призначення виконувачів обов'язків (в.о.) посадових осіб ОСС, призначених СУД відповідно до пункту 3.2 Положення про ОСС.

    5.6.2. Рішення про затвердження призначення в.о. приймається простою більшістю голосів від загального складу делегатів КСУ на найближчому засіданні після такого призначення.

    5.6.3. У разі відхилення КСУ призначення в.о., повноваження такого в.о. припиняються з моменту прийняття відповідного рішення КСУ, а питання про призначення нового в.о. розглядається у встановленому порядку.

    5.6.4. До моменту затвердження КСУ призначення в.о. має тимчасові повноваження відповідно до рішення СУД.